\chapter{Introduction}

\section{Motivation}

Large language models have shown strong ability in generating natural-language mathematical reasoning and formal proof search \citep{polu2020gptf,yang2023leandojo,xin2024deepseekproverv15}. 
However, a natural-language proof can hide many ambiguities: a step may omit assumptions, use an unstated lemma, rely on an informal transformation, or leave a gap that is accepted by the reader but not logically justified.
Formal proof assistants address this problem by requiring every proof step to be checked in a precise logical system. Lean 4 is one such proof assistant \citep{demoura2021lean4}. In Lean, mathematical statements and proofs are represented in a dependent type theory, and a theorem is accepted only when Lean can elaborate the proof and check that it has the required type. This makes Lean a useful setting for evaluating whether a generated mathematical proof is not only plausible, but formally correct.

Recent LLM-assisted proving systems therefore use large language models to propose Lean tactics, proof fragments, or complete proofs. 
However, the LLM is not the final judge of correctness. Lean remains the verifier: generated proof code must follow the required syntax, use available declarations, match the expected types, and close all goals.
This turns theorem proving into a search problem: the LLM proposes candidate proof steps, Lean verifies them, and the search procedure decides which proof states to explore next. 
At the same time, strict verification also makes proof generation fragile. A small local error, such as a missing delimiter, an unavailable theorem name, or an expression with the wrong type, can invalidate the whole proof attempt.

Proving methods can be roughly grouped into three directions: 

The first direction is whole-proof generation, where the model produces a complete proof in one pass. This is simple and can benefit from the model's long-form reasoning ability, but it is brittle: one local syntax, a declaration that is not available in the current context, or type error can cause the whole proof attempt to fail.

The second direction is tactic-level search. In Lean, a tactic transforms the current proof state into one or more new proof states. The model proposes local tactics, Lean executes and checks them, and the search procedure continues to explore the new proof states. Systems such as ReProver \cite{yang2023leandojo} combine tactic generation with best-first search. This gives backtracking and partial verification, but it also creates a large search space. Many different tactic sequences may be applicable at each proof state, each tactic may require different arguments or library lemmas, and every successful tactic can generate additional proof states that must be explored later. As proofs become longer, the search tree can grow combinatorially, requiring a large number of model samples and Lean executions before a useful proof path is found.

The third direction is proof decomposition. Instead of asking the model to solve the theorem directly or only propose the next tactic, the system asks the model to introduce useful intermediate steps or subgoals. This direction is motivated by the way long mathematical proofs are usually organized: a difficult theorem is reduced to smaller goals, and the final proof is assembled from those goals. Recent work such as ProD \cite{ProD} and Seed-Prover \cite{seedProver} has explored lemma-based or sketch-based decomposition for formal theorem proving.
Proof decomposition can reduce the difficulty of long-horizon reasoning and expose useful proof structure. However, decomposition also introduces new challenges. A generated subgoal may be mathematically true but irrelevant to the final proof. Large decompositions can introduce many additional subgoals, causing the search space to grow quickly. This also introduces a difficult credit assignment problem. A subgoal may be locally solvable but still provide little progress toward the original theorem. The usefulness of a decomposition is often visible only after several additional proof steps or after the final proof structure is completed.

These challenges make decomposition both promising and difficult: the system must not only generate useful subgoals, but also decide which decompositions are worth exploring, how to verify them, and how to connect them back to the final proof.

\section{Proposed Approach}

We propose GammaZero, a hierarchical proof-search framework for Lean 4. Instead of treating generated proofs as purely correct or incorrect, the system preserves useful local structure from partially successful proof attempts and uses it during search. The model proposes local proofs and decompositions, Lean verifies them, and the search graph records which fragments are valid, which subgoals are created, and which dependencies are needed for the final proof. 

We separate proof search into two complementary modes: directly solving the current goal and introducing proof structure with child subgoals. This leads to two corresponding action types:
\begin{itemize}
  \item \textbf{Local proof actions} are complete local proof attempts intended to discharge the current goal. A local proof action is accepted only if Lean reports no remaining goals for the current node.
  \item \textbf{Skeleton actions} are decompositions with proof structure, such as \texttt{have}, \texttt{suffices}, \texttt{cases}, or \texttt{induction}. A skeleton action may create child subgoals together with a partial proof structure describing how those subgoals are combined to solve the final theorem.
\end{itemize}

This separation turns proof search into an AND/OR graph. OR nodes represent proof states, while action nodes represent candidate proof steps. Each action forms an AND node with two types: a local proof action usually attempts to close the current goal directly and a skeleton action may create several child subgoals. A skeleton succeeds only when all of its child subgoals are solved.

To make failed generations useful for later learning, we introduce an AST-guided Sorrifier that extends placeholder-based recovery approaches such as APOLLO \cite{apollo}. The module localizes proof body failures and replaces the smallest recoverable failing region with \texttt{sorry}, producing a compilable partial script with placeholders. This patched script is not treated as a proof. It is used only to measure syntactic and structural validity.

The resulting system records how generated proofs evolve during search: which proof states are explored, which decompositions create useful subgoals, which fragments remain valid after recovery, and which subgoals contribute to the final proof.
The model proposes; Lean verifies; the search graph records alternatives and dependencies; the
Sorrifier recovers local signal from invalid code; dependency analysis measures whether skeleton
subgoals matter to the final proof; and final stitching reconstructs a proof candidate that Lean checks
again. Only complete Lean verification determines final correctness, but partially valid generations can still guide search and provide intermediate reward signals.

\section{Contributions}

This thesis makes the following contributions:

\begin{itemize}
    \item A hierarchical proof-search framework for Lean 4 that separates local proof generation from proof decomposition inside an AND/OR graph.

    \item A best-first search algorithm over Lean proof states using a priority queue and heuristic scoring.
    
    \item A scaffold method for solving subgoals inside the original parent proof context instead of treating them as independent theorems.
    
    \item A data synthesis pipeline that uses abstract syntax trees to recover partially valid Lean code, analyzes dependencies through Lean expression trees to evaluate the proof quality and converts proof search into structured trajectories with dense reward signals for future reinforcement learning.
\end{itemize}

\section{Thesis Organization}

Section \ref{sec:background} summarizes the Lean and LLM theorem-proving background needed for the proposed method. Section \ref{sec:literature} reviews prior theorem-proving systems and benchmarks. Section \ref{sec:problem-formulation} defines the proof state, action space, AND/OR objective, and main reward challenges. Section \ref{sec:framework} presents the end-to-end system pipeline. Section \ref{sec:Sorrifier} describes AST-guided recovery. Section \ref{sec:reward} defines dense reward and credit assignment. Section \ref{sec:experiments} reserves the experimental plan. Section \ref{sec:conclusion} concludes with limitations and next steps.
